\chapter{Arithmetic, Linear Algebra, Galerkin, and Splines}
\label{chap:arithmetic-galerkin-splines}

With a computer in hand, the construction can stop building tools and start using them. This chapter
does the mathematics the rest of the book will run on: ordinary arithmetic recovered as operations on
Facts, the linear algebra of bases and orthogonality, and then the two numerical methods that the
whole variational core depends on --- the Galerkin method, which solves problems by projecting them
into a finite frame, and splines, which stitch local pieces into a smooth whole. The chapter ends by
noticing something that the next several chapters will live on: a spline is already a variational
object, its joints first variations and its smooth fillers second variations. The calculus of
variations does not have to be imported. It is sitting inside the spline, waiting to be named.

\section{Values and addition}
\label{se:values-and-addition}

The first thing arithmetic needs is a value: a magnitude a measurement can carry, read off the
construction rather than assumed. The construction builds values, gives each a magnitude, and provides
the operation that compares two of them --- a decidable question of whether one magnitude overwhelms
another, which is all ``greater than'' has ever meant to a finite device. On top of values it builds
addition, and it builds it in a way that connects back to the fourth chapter: a sum is an accumulation,
the same act that turned a sequence of velocities into an area. To add is to lay one magnitude
alongside another and total the strip, so addition and integration are not two operations but one seen
at two scales --- a finite sum of finite pieces, with the integral as the limit the pieces only ever
approach.

Two small points make the arithmetic concrete. A magnitude comparison is not a fixed lookup but a
finite act: to ask whether one value overwhelms another, the construction walks their representations
from the high end down, the same bisection the earlier chapters used, and stops the moment one is
decided larger or the two prove equal to within resolution. And addition, read as area, is more than a
metaphor. Take three measured strips of widths two, three, and five; their sum, ten, is the area of the
region they tile end to end, and the same total would be read by integrating a step function that holds
each height across its width. The construction does not distinguish the two procedures, because for it
they are one: a finite total of finite pieces. What changes as the pieces shrink is only how closely
the staircase approximates a smooth ramp --- the limit the second chapter taught it to approach and
never to claim.

Repetition does something to values that no single value contains, and the construction borrows the
observation from Gauss. When the same measurement is added up again and again under a fixed way of
recording it, the totals do not scatter arbitrarily; they cluster into a stable bell-shaped
distribution with a definite center and a definite spread. The experiment under Gauss's name is careful
about why. The bell curve is not a law of nature smuggled in, and it does not assume the underlying
thing is smooth; it arises from the structure of repetition itself, the way independent finite errors
accumulate. Gauss introduced the normal curve exactly this way --- as the practical shape of error from
repeated observation with a fixed instrument, justified operationally rather than metaphysically --- and
the construction adopts it on the same terms. A mean and a variance are what repeated addition leaves
behind; they are summary Facts about a pile of trials, not properties of a world assumed in advance.

\section{Multiplication and orthogonality}
\label{se:multiplication-and-orthogonality}

If addition lays magnitudes end to end, multiplication scales them: it stretches or shrinks one
quantity by another, and the construction reads it as a spatial operation, a change of size rather than
a change of position. Two features of this scaling matter for what follows. The first is the sign flip.
Multiplying by a negative reverses orientation --- it turns a magnitude into its opposite --- and the
construction reads that reversal physically, as the flip between matter and antimatter, the same
sign-as-truth-value that the very first chapter built into its numbers. A value and its negative are
the same magnitude carried with opposite orientation, and the operation that exchanges them is
multiplication by minus one.

The second feature is orthogonality, and the construction gives it a reading that will recur for the
rest of the book. Two directions are orthogonal when scaling along one contributes nothing to the
other --- when their product cancels to zero. The construction reads this cancellation as destructive
interference: orthogonal components do not reinforce or diminish each other because they have no
overlap to share, and a measurement along one is blind to the other. This is the linear-algebraic
shadow of the superposition of the previous chapter --- where a blend of two states held both at once,
orthogonality is the condition under which two states can be held without contaminating each other ---
and it is the property that makes a basis useful: a set of mutually orthogonal directions is a set of
measurements that do not interfere, so each can be read cleanly without disturbing the rest.

A concrete pair fixes the idea. Two directions at right angles --- a step purely east and a step purely
north --- have a product that vanishes: moving east changes the eastward reading and leaves the
northward reading exactly as it was. That vanishing product is the whole content of orthogonality, and
it is why a surveyor can measure east and north separately and trust that neither contaminates the
other. Tilt the two directions toward each other and the product stops vanishing; now a step along one
leaks into the reading of the other, and the two measurements interfere. The construction's whole
strategy for handling many quantities at once is to choose directions whose products vanish in pairs,
so that each can be read in isolation --- and when no such directions are handed to it, to manufacture
them, straightening a tilted set into a right-angled one by subtracting off the overlaps. That
straightening is the orthogonalization every linear method begins with, and it is destructive
interference put to work: cancel the shared part on purpose, and what remains is clean.

\section{A basis and a small machine}
\label{se:a-basis-and-a-small-machine}

Linear algebra, stripped down, is the study of how to write any quantity as a combination of a chosen
few. The construction builds a basis --- a collection of independent directions --- and with it the two
structures every linear method needs: the span, the set of quantities the basis can reach by funging
its directions together in weighted combination, and the null space, the set of quantities that the
relevant transformation sends to nothing. To resolve a quantity in a basis is to tange its components,
selecting how much of each independent direction it carries, and the cleanness of that reading is
exactly the orthogonality of the previous section.

The construction cannot resist making this concrete as a tiny programmable machine, and the joke
carries a real idea. Picture the smallest interpreter that can branch and call: it has a jump, which
sends control somewhere else and abandons where it was, and a call, which goes somewhere else and
remembers to return. The construction reads the jump as scattering --- an interaction that sends a
state off in a new direction with no memory of the old --- and the call as a propagator, the operator
that carries a state forward and brings it back, the bookkeeping of how influence travels and returns.
Underneath both is the eigenvector reading: among all the directions a transformation can act on, the
eigenvectors are the ones it merely scales without turning, the directions a measurement can load and
read back unrotated. An eigenvector is a direction that survives a process intact, which is why it is
the natural thing for a basis to be built from --- a basis of eigenvectors is a set of measurements
each of which the process leaves pointing where it found it.

A worked resolution shows the machinery turning. Suppose a quantity is to be written in a basis of three
orthogonal directions. The construction asks, of each direction in turn, how much of it the quantity
carries --- a single projection, a product that does not vanish --- and records the three answers as the
quantity's components. Reassembled, the three weighted directions return the original exactly, because
the basis spanned the space and the orthogonality kept the questions from interfering. Now suppose the
basis is chosen from the eigenvectors of the process under study. Then the process, applied to the
quantity, does nothing to the directions and only rescales the components --- each eigenvector stretched
by its own factor and turned not at all. A problem that looked like a tangle of coupled motions becomes
three independent scalings, one per eigenvector, and that decoupling is the reason eigenvectors are
worth the trouble of finding: in the right basis, a hard process falls apart into easy ones.

\section{Galerkin and splines}
\label{se:galerkin-and-splines}

Now the chapter reaches the methods the rest of the book runs on. A polynomial is the construction's
flexible shape: a finite combination of basis powers that can be bent to approximate almost anything
over a small enough region. The Galerkin method is the way the construction solves an equation using
them, and its defining virtue is that it is derivative-free. Rather than differentiating --- which a
finite device cannot do exactly, only approximate --- it demands that the leftover error, the residue
of a candidate solution, be orthogonal to every polynomial in a chosen finite basis. If the residue
has no component along any direction the basis can see, then within that frame the candidate is the
solution, and the whole problem reduces to the orthogonality of the previous section applied to the
error. Differentiation is replaced by projection, and projection is just funging the residue into the
part the basis cannot detect.

Made concrete, the method is disarmingly simple. Suppose the construction wants to approximate some
target behavior by a polynomial drawn from the finite basis of a constant, a linear term, and a
quadratic. It writes its candidate as an unknown weighted sum of the three, computes the residue --- the
amount by which the candidate fails to satisfy the equation it is meant to solve --- and then imposes
exactly three conditions: that the residue have no overlap with the constant, none with the linear
term, and none with the quadratic. Three orthogonality conditions, three unknown weights; the
construction solves the small finite system and reads off the best approximation the three-term frame
can give. No derivative was ever taken. The error was not minimized by calculus but annihilated, in
each of the basis directions, by projection --- which is why the method survives in a world where exact
differentiation is unavailable but a finite inner product is always at hand.

The phrase ``finite basis'' is load-bearing, and the construction is blunt about its cost. The method
lives in a single finite frame --- one bounded stack of working memory, which the construction nicknames
to keep it honest about the fact that it is singular and only so big, and which a few keystrokes can
overflow. Everything Galerkin does, it must do within that one frame; it cannot recurse without limit or
hold an infinite basis. The check that the method has actually converged is correspondingly finite: a
transform is finite, the construction says, when it settles --- when it reaches a state the next step
leaves unchanged, an eigenvalue fixed point where the approximating has stopped, or when its residue
stops carrying new strain from one step to the next. A transform that would inject unbounded energy ---
moving from a settled ground state back into an excited one --- is ruled out, not as an error but as a
divergence the finite frame cannot represent. Convergence, here, is the statement that the method came
to rest inside the one frame it was given.

The convergence test deserves to be made precise, because it is where the method admits it can fail.
After each refinement the construction compares the new approximation with the one before it and asks a
decidable question: did anything of substance change? If the candidate has reached a state the next step
leaves fixed --- an eigenvalue fixed point, where projecting again returns the same polynomial --- the
method has converged and stops. If the residue keeps carrying the same strain forward without shrinking,
the method is stalled; and if a step would inject more than the frame can hold, driving a settled
solution back into an excited one, the construction rules it out as a divergence rather than chasing it.
There is no promise, issued in advance, that a given problem will settle inside the one frame --- that is
the same halting honesty the fourth chapter insisted on, now wearing the clothes of numerical analysis.
The method is trustworthy exactly when it comes to rest, and it tells you whether it has.

A spline is what you get when you stitch such local polynomial pieces together, and it is the
chapter's payoff because it is already the object the calculus of variations studies. A spline is built
from three things, and their names in the construction are exact. There is the observation, the bare
ground state, a recorded value with no work done on it. There is the knot --- a localized place where
two polynomial pieces meet --- and the construction marks it as the \emph{first variation}: a knot is
where the curve is allowed to bend, the place a small change in the path first registers. And there is
the interpolant, the smooth filler that joins the pieces across a knot, which the construction marks as
the \emph{second variation}, the weak derivative, the curvature that the bending leaves behind. A
spline, in other words, is not merely a smooth curve through data; it is a record of first and second
variations, a structure whose joints are the first-order freedom to vary and whose fillers are the
second-order response. The whole apparatus of stationary paths and minimal actions that the next
chapters build is, at bottom, the discipline of choosing those knots and interpolants well.

The experiment that names what this is for is the harmonic oscillator. Read informationally --- with no
metric and no assumed dynamics --- it is a conjugate pair: one variable recording the current
distinguishable state, its partner recording the rate at which that distinguishability is expected to
change, the two trading off under a reciprocity that makes the system oscillate. The quantity that
governs such a pair is a quadratic form, and a quadratic form is exactly a second variation. The
harmonic oscillator is the simplest place the second variation appears as the thing that runs the
dynamics, and it is the construction's signpost that the spline's interpolant --- its second variation
--- is the object the whole middle of the book will turn on.

\section{What is demonstrated, and what is assumed}
\label{se:demonstrated-assumed-ch7}

Following the experimental library's discipline, we separate the built from the named. This chapter
built a great deal, and most of it is solid in the plainest way: it is ordinary mathematics, recovered
finitely.

What is demonstrated is arithmetic and linear algebra over the construction's own numbers --- values
with a decidable magnitude order, addition as finite accumulation, multiplication as scaling with a
sign flip, a basis with its span and null space, and the eigenvector reading --- together with two
numerical methods: a derivative-free Galerkin projection that solves by making a residue orthogonal to
a finite basis, and a spline built from observations, knots, and interpolants. Every piece is finite
and decidable, and the Galerkin method's reliance on a single bounded frame is stated as a genuine
limitation, not hidden.

\begin{quote}
\emph{A note on the labels.} The physical readings here are bridges. Calling the sign flip a
\emph{matter--antimatter} exchange, orthogonality \emph{destructive interference}, a jump and a call
\emph{scattering} and a \emph{propagator}, and the spline's interpolant a \emph{second variation} are
interpretations laid over finite linear algebra. The mathematics is earned outright; the physical names
are the analogy the construction is betting on. The single identification that is \emph{not} an analogy
is internal and will be cashed in over the next several chapters: that a spline's knot is a first
variation and its interpolant a second. That one is a fact about the construction, and the calculus of
variations is what happens when it is taken seriously.
\end{quote}

What is assumed is modest here and mostly inherited: that the finite Galerkin frame is large enough for
the problems put to it, and that the physical names above are the right ones for the linear-algebraic
structures underneath. With arithmetic, linear algebra, and the spline in hand --- and with the second
variation already glinting inside the spline's interpolant --- the construction has everything it needs
to turn to the question the whole next part is built around: among all the paths a measurement could
take, which one does it actually take, and what holds steady when it is nudged.
